
\assignmenttitle
    {LAB ASSIGNMENT 2}
    {POINTERS, STRUCTURES, FUNCTIONS}

\aim{
    To write C++ programs utilizing concepts of pointers, structures and functions.
}

\programs

\program{
    1. A phone number, such as (212) 767-8900, can be thought of as having three parts: the area code (212), the exchange (767), and the number (8900). Write a program that uses a structure to store these three parts of a phone number separately. Call the structure phone. Create two structure variables of type phone. Initialize one, and have the user input a number for the other one. Then display both numbers. The interchange might look like this:\\[4pt]
    Enter your area code, exchange, and number: 415 555 1212\\[4pt]
    My number is (212) 767-8900\\[4pt]
    Your number is (415) 555-1212
}

\codelabel

\begin{code}
#include <iostream>
using namespace std;

struct phone {
    int area, exchange, number;
};

int main() {
    phone my = {212, 767, 8900}, your;
    cout << "Enter your area code, exchange and number: ";
    cin >> your.area >> your.exchange >> your.number;
    cout << "My number is (" << my.area << ") " << my.exchange << "-" << my.number << endl;
    cout << "Your number is (" << your.area << ") " << your.exchange << "-" << your.number << endl;
    return 0;
}
\end{code}

\outputlabel
\outputimage[0.50\linewidth]{images/OOP/image11.png}

\program{
    2. A point on the two-dimensional plane can be represented by two numbers: an x coordinate and a y coordinate. For example, (4,5) represents a point 4 units to the right of the vertical axis, and 5 units up from the horizontal axis. The sum of two points can be defined as a new point whose x coordinate is the sum of the x coordinates of the two points, and whose y coordinate is the sum of the y coordinates.\\[4pt]
    Write a program that uses a structure called point to model a point. Define three points, and have the user input values to two of them. Then set the third point equal to the sum of the other two, and display the value of the new point. Interaction with the program might look like this:\\[4pt]
    Enter coordinates for p1: 3 4\\[4pt]
    Enter coordinates for p2: 5 7\\[4pt]
    Coordinates of p1+p2 are: 8, 11
}

\codelabel

\begin{code}
#include <iostream>
using namespace std;

struct point {
    int x, y;
};

int main() {
    point p1, p2, p3;
    cout << "Enter coordinates for p1: ";
    cin >> p1.x >> p1.y;
    cout << "Enter coordinates for p2: ";
    cin >> p2.x >> p2.y;
    p3.x = p1.x + p2.x;
    p3.y = p1.y + p2.y;
    cout << "Coordinates of p1+p2 are: " << p3.x << ", " << p3.y;
    return 0;
}
\end{code}

\outputlabel
\outputimage[0.50\linewidth]{images/OOP/image12.png}

\program{
    3. Create a structure called time. Its three members, all type int, should be called hours, minutes, and seconds. Write a program that prompts the user to enter a time value in hours, minutes, and seconds. This can be in 12:59:59 format, or each number can be entered at a separate prompt (``Enter hours:'', and so forth). The program should then store the time in a variable of type struct time, and finally print out the total number of seconds represented by this time value.
}

\codelabel

\begin{code}
#include <iostream>
using namespace std;

struct time {
    int hours, minutes, seconds;
};

int main() {
    struct time t;
    cout << "Enter hours: ";
    cin >> t.hours;
    cout << "Enter minutes: ";
    cin >> t.minutes;
    cout << "Enter seconds: ";
    cin >> t.seconds;
    cout << "Total seconds = " << t.hours * 3600 + t.minutes * 60 + t.seconds;
    return 0;
}
\end{code}

\outputlabel
\outputimage[0.50\linewidth]{images/OOP/image13.png}

\program{
    4. Implement a function Swap (int *a, int *b) that swaps the values of two integers using pointers. Demonstrate its use in a main function that accepts two integers and prints them before and after swapping.
}

\codelabel

\begin{code}
#include <iostream>
using namespace std;

void Swap(int *a, int *b) {
    int temp = *a;
    *a = *b;
    *b = temp;
}

int main() {
    int x, y;
    cout << "Enter two numbers: ";
    cin >> x >> y;
    cout << "Before Swap: " << x << " " << y << endl;
    Swap(&x, &y);
    cout << "After Swap: " << x << " " << y;
    return 0;
}
\end{code}

\outputlabel
\outputimage[0.50\linewidth]{images/OOP/image14.png}

\program{
    5. Define inline functions for basic arithmetic operations (add, subtract, multiply, divide) and use them in a menu-driven calculator program that performs operations based on user input.
}

\codelabel

\begin{code}
#include <iostream>
using namespace std;

inline int add(int a, int b) { return a + b; }
inline int sub(int a, int b) { return a - b; }
inline int mul(int a, int b) { return a * b; }
inline float divi(int a, int b) { return (float)a / b; }

int main() {
    int a, b, ch;
    cout << "Enter two numbers: ";
    cin >> a >> b;
    cout << "1.Add\n2.Subtract\n3.Multiply\n4.Divide\n";
    cin >> ch;
    if (ch == 1) cout << add(a, b);
    else if (ch == 2) cout << sub(a, b);
    else if (ch == 3) cout << mul(a, b);
    else if (ch == 4) cout << divi(a, b);
    else cout << "Invalid Choice";
    return 0;
}
\end{code}

\outputlabel
\outputimage[0.50\linewidth]{images/OOP/image15.png}

\program{
    6. Write a recursive function that accepts a pointer to an array and computes the sum of elements. Extend this to find the maximum value recursively, using pointer arithmetic.
}

\codelabel

\begin{code}
#include <iostream>
using namespace std;

int sum(int *arr, int n) {
    if (n == 0) return 0;
    return *arr + sum(arr + 1, n - 1);
}

int maximum(int *arr, int n) {
    if (n == 1) return *arr;
    int m = maximum(arr + 1, n - 1);
    return (*arr > m) ? *arr : m;
}

int main() {
    int n;
    cout << "Enter size: ";
    cin >> n;
    int arr[n];
    for (int i = 0; i < n; i++)
        cin >> arr[i];
    cout << "Sum = " << sum(arr, n) << endl;
    cout << "Maximum = " << maximum(arr, n);
    return 0;
}
\end{code}

\outputlabel
\outputimage[0.50\linewidth]{images/OOP/image16.png}

\program{
    7. Start with an array of pointers to strings representing the days of the week. Provide functions to sort the strings into alphabetical order. Sort the pointers to the strings, not the actual strings.
}

\codelabel

\begin{code}
#include <iostream>
#include <cstring>
using namespace std;

int main() {
    char *days[] = {
        "Sunday", "Monday", "Tuesday", "Wednesday",
        "Thursday", "Friday", "Saturday"
    };

    for (int i = 0; i < 7; i++) {
        for (int j = i + 1; j < 7; j++) {
            if (strcmp(days[i], days[j]) > 0) {
                char *temp = days[i];
                days[i] = days[j];
                days[j] = temp;
            }
        }
    }

    cout << "Sorted Days:\n";
    for (int i = 0; i < 7; i++)
        cout << days[i] << endl;
    return 0;
}
\end{code}

\outputlabel
\outputimage[0.50\linewidth]{images/OOP/image17.png}

\program{
    8. Simulate how arrays are stored in memory by writing a program that allocates an array dynamically and then displays both the content and the address of each element. Use pointer arithmetic to traverse the array and show the difference in address steps.
}

\codelabel

\begin{code}
#include <iostream>
using namespace std;

int main() {
    int n;
    cout << "Enter size: ";
    cin >> n;
    int *arr = new int[n];

    for (int i = 0; i < n; i++) {
        cout << "Enter element: ";
        cin >> *(arr + i);
    }

    cout << "\nValue    Address\n";
    for (int i = 0; i < n; i++)
        cout << *(arr + i) << "       " << (arr + i) << endl;

    delete[] arr;
    return 0;
}
\end{code}

\outputlabel
\outputimage[0.50\linewidth]{images/OOP/image18.png}

\program{
    9. Simulate a memory manager that keeps track of memory blocks (e.g., like malloc/free) using an array and pointers. Allow users to allocate, deallocate, and compact memory. Use pointer arithmetic to manage and traverse the memory layout.
}

\codelabel

\begin{code}
#include <iostream>
using namespace std;

int main() {
    int memory[10] = {0}, choice, pos;

    while (1) {
        cout << "\n1.Allocate\n2.Free\n3.Display\n4.Exit\n";
        cin >> choice;
        if (choice == 1) {
            cout << "Enter position to (0-9): ";
            cin >> pos;
            memory[pos] = 1;
        } else if (choice == 2) {
            cout << "Enter position: ";
            cin >> pos;
            memory[pos] = 0;
        } else if (choice == 3) {
            for (int *p = memory; p < memory + 10; p++)
                cout << *p << " ";
            cout << endl;
        } else
            break;
    }
    return 0;
}
\end{code}

\outputlabel
\outputimage[0.50\linewidth]{images/OOP/image19.png}
\outputimage[0.50\linewidth]{images/OOP/image20.png}
\outputimage[0.50\linewidth]{images/OOP/image21.png}
\outputimage[0.50\linewidth]{images/OOP/image22.png}
\outputimage[0.50\linewidth]{images/OOP/image23.png}

